% !TeX program = xelatex
% Standalone simulation report.  Mathematical exposition is handwritten;
% every numerical table and plotted coordinate is produced by the shared
% Dask pipeline.  The original numerical experiment and evidence are retained;
% the mathematical exposition has a separate proof-revision manifest.
% Shared XeLaTeX report style. Charts consume Dask-produced CSVs via PGFPlots.
\documentclass[a4paper,fontsize=11pt,parskip=half,headings=small]{scrartcl}
\usepackage[margin=18mm,includeheadfoot,headheight=14pt,headsep=6mm,footskip=9mm]{geometry}
\usepackage{fontspec}
\usepackage{mathtools,amsthm}
\usepackage{unicode-math}
\setmainfont{TeX Gyre Pagella}
\setsansfont{TeX Gyre Heros}[Scale=MatchLowercase]
\setmonofont{Latin Modern Mono}
\setmathfont{TeX Gyre Pagella Math}
\usepackage[protrusion=true,expansion=false]{microtype}
\usepackage{xcolor,xurl,booktabs,longtable,array,needspace,placeins}
\usepackage[numbers,sort&compress]{natbib}
\usepackage{pgfplots}
\pgfplotsset{compat=1.18,width=.92\linewidth,height=6.0cm,
 every axis/.append style={grid=major,grid style={gray!20},
 tick label style={font=\small},label style={font=\small},
 legend style={font=\small,draw=none,fill=white}}}
\usepgfplotslibrary{groupplots}
\usepackage[unicode,colorlinks=true,linkcolor=blue!50!black,citecolor=blue!50!black,urlcolor=blue!50!black]{hyperref}
\usepackage{bookmark,scrlayer-scrpage}
\setlength{\emergencystretch}{1.5em}
\setkomafont{disposition}{\normalfont\sffamily\bfseries}
\setkomafont{pageheadfoot}{\normalfont\sffamily\small}
\clearpairofpagestyles
\ihead{Simulation methods and results}
\ohead{Dask experiment report}
\cfoot*{\pagemark}
\pagestyle{scrheadings}
\newtheorem{theorem}{Theorem}[section]
\newtheorem{proposition}[theorem]{Proposition}
\newtheorem{lemma}[theorem]{Lemma}
\newtheorem{corollary}[theorem]{Corollary}
\DeclareMathOperator{\E}{E}
\DeclareMathOperator{\Var}{Var}
\DeclareMathOperator{\Cov}{Cov}
\DeclareMathOperator{\SD}{SD}

\hypersetup{pdftitle={Sequential Conditional Means: Simulation Methods and Results},
  pdfsubject={Dask experiments with population and conditional actual-score means}}
\title{Sequential Conditional Means:\\Simulation Methods and Results}
\subtitle{One, two or three initial IQ-150 scores, with every initial score included}
\author{\ProjectAuthor\\{\normalsize\ProjectRoles}}
\date{15 November 2026}

\begin{document}
\frontmatter
% Computed by the Dask report task; display rounding only.
\newcommand{\SimTrials}{1\,000\,000}
\newcommand{\SimReplicates}{4}
\newcommand{\SimPooledTrials}{4\,000\,000}
\newcommand{\SimChunk}{50\,000}
\newcommand{\SimMaxMeanSE}{2.549}
\newcommand{\SimMaxVarianceRatio}{0.002317}
\newcommand{\SimRareMin}{0}
\newcommand{\SimRareMax}{0}
\newcommand{\SimMaxEventSE}{2.630}
\newcommand{\SimMaxSDRelativeError}{0.001158}

\maketitle
\thispagestyle{plain.scrheadings}

\begin{abstract}
This report explains the simulations accompanying the sequential
conditional-means analysis. Standardized tests have population mean zero,
standard deviation one and pairwise correlation \(c=0.6\). One, two or
three observed starting scores are fixed at IQ 150, followed by one, two
or ten actual future scores. We simulate whole population batteries and
conditional future vectors as separate experiments, compare their means
and standard deviations with the Gaussian formulas, and sample the
Gaussian laws with the calculated infinite-count parameters. Each case
uses \SimReplicates{} seeded replicates of \SimTrials{} observations,
pooled to \SimPooledTrials{} observations. Dask builds and executes the
array, reduction and report-data graphs; XeLaTeX/PGFPlots draws figures
from their generated CSVs. The key distinction is between variation in
actual outcomes and the zero standard deviation of a fixed list of
predictions. Agreement is evidence about this implementation of the
assumed model, not empirical validation of that model for IQ data.
\end{abstract}

\noindent\textbf{AI disclosure.} OpenAI Codex and collaborating AI agents
prepared this report. The full disclosure appears in the final appendix,
\Cref{sec:ai-disclosure}.

\noindent\textbf{Proof and evidence convention.} The Gaussian assumptions,
the reporting scale and the experiment design are stated inputs. Every
deductive target and uncertainty statement has a proof in this report;
\Cref{sec:seq-sim-proof-map} links the claim groups to those proofs
and the established theorems they use. Observed numerical values and
software-check counts are finite reproduction records, with their scope
specified in \Cref{seqp:record-scope}.

\tableofcontents
\mainmatter

\chapter{The questions answered by the experiments}
\label{sec:seq-sim-questions}

The companion analysis studies an inclusive average: every observed
starting score contributes to the numerator and denominator. Write
\(r\in\{1,2,3\}\) for the number of initial scores, \(k\in\{1,2,10\}\)
for the number of later scores and \(N=r+k\) for the total. All initial
scores have the same standardized value
\[
 a=\frac{150-100}{15}=\frac{10}{3},\qquad Q=100+15X.
\]
The reporting transformation changes units; it does not change the
underlying requirement that each standardized test has population mean
zero and variance one.

There are three quantities to keep separate. The \emph{population actual
mean} comes from a newly drawn whole battery, including newly drawn
initial scores. Its IQ centre is 100. The \emph{conditional actual mean}
comes from future outcomes after the initial \(r\) scores have been
observed at 150; its centre is the conditional mean calculated below.
Finally, the \emph{recursively filled mean} replaces every future outcome
by its prediction. Once \(r,k,c,a\) are fixed, that is a number rather
than a new random outcome, and its standard deviation is zero.

The simulations ask whether separately generated whole-battery and
conditional outcomes reproduce the appropriate theoretical centres and
spreads. They also ask how those spreads change with the number of
future scores, and whether the infinite-count parameters are consistent
with the finite formulas. They do not regard a narrow Monte Carlo
confidence interval for a simulated centre as a narrow prediction
interval for a person's actual future score mean.

The three means are derived in Proposition~\ref{seqp:recursion}.
Proposition~\ref{seqp:variance-decomposition} separates their sources of
variance. The distinction between outcome spread and estimation uncertainty
follows from Proposition~\ref{sm:normal-sample} and
Corollary~\ref{sm:sd-interval}.

\chapter{Referenced foundations and full working proofs}
\label{sec:proof-foundations}

The results below supply the mathematical foundations used in this report.
The explicitly assumed inputs are finite second moments and, where stated,
a jointly Gaussian law and independent observations. An assumption is not a
consequence of a simulation. Definitions, source quotations, historical
records, rounded numerical evaluations and recommendations have the status
specified in the claim map; they are not additional mathematical theorems.
The proofs reduce the deductions to linearity of integration, the spectral
theorem, the Gaussian characteristic function and conditional expectation.
Source locators refer to the linked author versions specified in the
bibliography, rather than silently mixing their pagination with print editions.

\noindent\textbf{Document-revision provenance.}
The current book-format revision is recorded in
\path{reporting/book-format/revision.json}, with current TeX/Bib sources
and PDF hashes. The earlier proof-text revision remains recorded in
\path{reporting/proof-revision/revision.json}.
The original \path{reporting/run.json} retains the numerical
simulation campaign, its original document hashes and numerical outputs.
The book revision checks and reuses the unchanged simulation-data files; it does
not rerun the experiment or claim a new native SPARK/Rocq proof.
The repository's subsequent shared-library refactor is separate from the
retained native campaigns. Their source-pinned results do not constitute
new GNATprove evidence for refactored Ada bodies; those changes require
their own matching campaign. The written mathematical proofs below do not
depend on that program-verification status.

\begin{proposition}[Finite moments and affine changes of scale]
\label{pf:moments}
Let a real random vector \(X\) have finite second moments, mean \(\mu\)
and covariance \(\Sigma\). Let \(Y\) be square-integrable, \(w\) a
deterministic vector of matching dimension, and \(b,s\) real constants. Then
\[
 \E(w^{\mathsf T}X)=w^{\mathsf T}\mu,\qquad
 \Var(w^{\mathsf T}X)=w^{\mathsf T}\Sigma w,\qquad
 \SD(b+sY)=|s|\SD(Y).
\]
A zero variance is equivalent to equality with the mean almost surely.
\end{proposition}
\begin{proof}
H\"older's inequality ensures integrability of all products
\((X_i-\mu_i)(X_j-\mu_j)\). Finite linearity of the integral then gives
the expectation formula and
\[
 \E\!\left[\Big(\sum_iw_i(X_i-\mu_i)\Big)^2\right]
 =\sum_{i,j}w_iw_j\E[(X_i-\mu_i)(X_j-\mu_j)]
 =w^{\mathsf T}\Sigma w.
\]
Replacing \(Y\) by \(b+sY\) multiplies its centered value by \(s\), so
its variance is \(s^2\Var(Y)\); taking the nonnegative square root proves
the SD formula. If \(\E[(Y-\E Y)^2]=0\), Markov's inequality applied to
\((Y-\E Y)^2\) makes \(P(|Y-\E Y|>\epsilon)=0\) for every
\(\epsilon>0\). The countable union over \(\epsilon=1/n\) proves the
almost-sure assertion. The reverse implication follows by integration.
These are applications of linearity, H\"older and the variance identities
in \citet[Section 1.6, especially Theorems 1.6.1, 1.6.3--1.6.4]{durrett2019}.
\end{proof}

\begin{proposition}[Covariance feasibility and Gaussian construction]
\label{pf:spectral}
A real symmetric matrix is a covariance matrix if and only if it is
positive semidefinite. For the equicorrelation matrix
\(R_n=(1-c)I_n+c\mathbf1_n\mathbf1_n^{\mathsf T}\), \(n\geq2\), this
is equivalent to \(-1/(n-1)\leq c\leq1\). Positive definiteness requires
both inequalities to be strict.
\end{proposition}
\begin{proof}
Proposition~\ref{pf:moments} gives
\(w^{\mathsf T}\Sigma w=\E[(w^{\mathsf T}(X-\mu))^2]\geq0\).
Conversely, the real spectral theorem diagonalizes a symmetric positive
semidefinite matrix as \(QDQ^{\mathsf T}\), with orthogonal \(Q\) and
nonnegative diagonal entries \(D\). For independent standard normal
coordinates \(Z\), the vector \(QD^{1/2}Z\) has covariance
\(QD^{1/2}I D^{1/2}Q^{\mathsf T}=\Sigma\) and is Gaussian, including
zero eigenvalues. The spectral theorem and positive square roots are
\citet[Theorems 7.29 and 7.38]{axler2024}; the Gaussian construction is
also \citet[Theorems 3.4.4--3.4.5]{gallager2013}.
For \(R_n\), the span of \(\mathbf1_n\) has eigenvalue
\(1+(n-1)c\), and its orthogonal complement has eigenvalue \(1-c\),
because \(\mathbf1_n^{\mathsf T}w=0\) there. Every vector decomposes
into these two orthogonal subspaces. Nonnegativity, or positivity, of
both eigenvalues gives precisely the stated inequalities. If one fixed
\(c\) must work for arbitrarily large \(n\), the lower bounds force
\(c\geq0\); the finite-dimensional bounds alone permit negative \(c\).
\end{proof}

\begin{proposition}[Gaussian closure and block independence]
\label{pf:gaussian}
An affine image of a jointly Gaussian vector is jointly Gaussian. Its law
is determined by its mean and covariance. Two blocks of a jointly Gaussian
vector are independent whenever their cross-covariance is zero, including
singular Gaussian vectors.
\end{proposition}
\begin{proof}
For a Gaussian vector with mean \(\mu\) and covariance \(\Sigma\), its
characteristic function is
\(\psi_X(t)=\exp(i t^{\mathsf T}\mu-t^{\mathsf T}\Sigma t/2)\).
Thus \(\psi_{AX+b}(t)=e^{it^{\mathsf T}b}\psi_X(A^{\mathsf T}t)\)
has the Gaussian form with mean \(A\mu+b\) and covariance
\(A\Sigma A^{\mathsf T}\). This also follows from the independent-normal
affine representation, without invertibility. If a two-block covariance
has zero off-diagonal blocks, the displayed characteristic function
factors as \(\psi_U(s)\psi_V(t)\). By Fubini this is the characteristic
function of the product of the two marginal measures. Uniqueness of
characteristic functions proves both independence and determination by
mean and covariance. The Gaussian formula and closure are
\citet[Theorems 3.3.1--3.3.2, equations (3.14) and (3.18)]{gallager2013};
uniqueness is \citet[Theorem 3.10.4]{durrett2019}.
\end{proof}

\begin{proposition}[Gaussian conditional distribution, including exact histories]
\label{pf:conditioning}
Suppose \((U,V)\) is jointly Gaussian with mean \((\mu_U,\mu_V)\),
covariance blocks \(A,B\) and \(C=\Cov(U,V)\), and \(B\) positive
definite. A version of the conditional law, for every \(v\), is Gaussian
with mean and covariance
\[
 \mu_U+CB^{-1}(v-\mu_V),\qquad A-CB^{-1}C^{\mathsf T}.
\]
The conditional-mean formula is the unique continuous version on the
whole conditioning space. The full joint covariance need not be invertible.
\end{proposition}
\begin{proof}
Set \(L=CB^{-1}\) and
\(W=U-\mu_U-L(V-\mu_V)\). Its cross-covariance with \(V\) is
\(C-LB=0\). Expanding its covariance gives
\(A-LC^{\mathsf T}-CL^{\mathsf T}+LBL^{\mathsf T}
=A-CB^{-1}C^{\mathsf T}\). Proposition~\ref{pf:gaussian} makes
\(W\) Gaussian and independent of \(V\). Consequently the kernel
\(K(v,H)=P(\mu_U+L(v-\mu_V)+W\in H)\) satisfies
\[
 P(U\in H,V\in D)=\int_D K(v,H)P_V(dv)
\]
by independence and Fubini, which is the defining conditional-law identity.
Its expectation is the displayed affine function. A positive-definite
Gaussian \(V\) has a strictly positive density everywhere. If two
continuous versions differed at any point, continuity would give a
neighborhood of positive probability on which they differ, contradicting
almost-sure uniqueness. Conditioning on an exact continuous history uses
this specified version, not division by the probability-zero event
\(\{V=v\}\). This residual proof is the Gaussian regression theorem
\citep[Theorems 3.5.1--3.5.2]{gallager2013}, proved here in the slightly
more general form allowing singular \((U,V)\); conditional versions are
defined in \citet[Section 4.1]{durrett2019}.
\end{proof}

\begin{proposition}[Normal marginal distributions do not imply joint normality]
\label{pf:normal-marginals}
Normal marginals and a specified Pearson correlation do not in general
determine Gaussian composite or joint-tail laws.
\end{proposition}
\begin{proof}
Let \(X\sim N(0,1)\) and independently let \(S\) equal \(1\) with
probability \((1+\rho)/2\) and \(-1\) otherwise, \(-1<\rho<1\).
Put \(Y=SX\). Symmetry gives \(Y\sim N(0,1)\), and independence gives
\(\E[XY]=\E[S]\E[X^2]=\rho\). Nevertheless
\(P(X+Y=0)=(1-\rho)/2>0\). A nondegenerate normal variable has a
density and no atoms; \(X+Y\) has variance \(2+2\rho>0\), so this
pair is not jointly Gaussian. For a concrete tail nonidentification at
\(\rho=0\), let \(t>0\) and \(p=P(X>t)\in(0,1/2)\). This pair
has \(P(X>t,Y>t)=p/2\), whereas an independent standard-normal pair,
with the same marginals and zero correlation, has intersection \(p^2\).
These differ since \(p<1/2\). The construction uses symmetry,
independence and the product-integration rule
\citep[Sections 1.7 and 2.1]{durrett2019}; it is a counterexample,
not an alternative assumption secretly used for the report's Gaussian tables.
\end{proof}

\begin{proposition}[Event arithmetic and conditional denominators]
\label{pf:events}
For events \(A,B\) with probabilities \(p_A,p_B\) and intersection
\(J\), their union has probability \(p_A+p_B-J\). If its conditioning
probability is positive, \(P(B\mid A)=J/p_A\). The correlation of the
nonconstant indicators is
\[
 \frac{J-p_Ap_B}{\sqrt{p_A(1-p_A)p_B(1-p_B)}}.
\]
\end{proposition}
\begin{proof}
The pointwise identities
\(1_{A\cup B}=1_A+1_B-1_{A\cap B}\) and
\(1_A1_B=1_{A\cap B}\) prove the union and covariance formulas by
linearity. Because \(1_A^2=1_A\),
\(\Var(1_A)=p_A-p_A^2\), and similarly for \(B\); divide their
covariance by the product of their positive SDs. The conditional formula
is precisely the event-conditional probability definition, requiring
\(p_A>0\). The four disjoint cells have probabilities
\(J,p_A-J,p_B-J,1-p_A-p_B+J\), obtained by the same identities.
These are consequences of finite additivity and conditional probability
\citep[Sections 1.1 and 4.1]{durrett2019}.
\end{proof}

\begin{proposition}[Tower property and population--conditional variance]
\label{pf:expectation}
For square-integrable \(Y\) and a sub-\(\sigma\)-algebra \(\mathcal G\),
\[
 \E\E[Y\mid\mathcal G]=\E Y,\qquad
 \Var(Y)=\E\Var(Y\mid\mathcal G)
              +\Var(\E[Y\mid\mathcal G]).
\]
\end{proposition}
\begin{proof}
The defining conditional-integral identity with the whole space gives
the first formula. Write \(M=\E[Y\mid\mathcal G]\). Conditional
Jensen \citep[Theorem 4.1.10]{durrett2019} ensures \(M\) is
square-integrable. In
\(Y-\E Y=(Y-M)+(M-\E Y)\), the cross term has expectation zero:
the pull-out and tower properties give
\(\E[(Y-M)(M-\E Y)]
=\E[(M-\E Y)\E(Y-M\mid\mathcal G)]=0\).
Expanding the square and integrating gives the variance identity. The
pull-out rule extends from bounded multipliers by truncation and
Cauchy--Schwarz. These conditional-expectation properties follow from
\citet[Section 4.1, especially Theorems 4.1.13--4.1.14]{durrett2019}.
In particular, fixing an observed starting score and varying it across
the population are different probability experiments, even when both
use the same marginal test SD.
\end{proof}


\chapter{The Gaussian model and its exact target values}
\label{sec:seq-sim-model}

\section{A model with a common factor}
The experiments assume the jointly Gaussian construction
\begin{equation}
 X_i=\sqrt c\,G+\sqrt{1-c}\,E_i,
 \label{eq:seq-sim-factor}
\end{equation}
where \(G,E_0,E_1,\ldots\) are independent standard normals. By linearity
and independence, each \(X_i\) has mean zero and variance
\(c+(1-c)=1\), and two distinct tests have covariance \(c\). Thus
\eqref{eq:seq-sim-factor} realizes the stipulated equicorrelation model.
Gaussian affine closure and independence of Gaussian blocks with zero
cross-covariance are established reference results
\citep[Chapter 3, especially Theorems 3.3.1--3.3.2]{gallager2013}.
The companion derivation proves their particular conditional-regression
consequences, including the choice of continuous conditional version.
For self-contained proofs here, see the Gaussian foundation in
Proposition~\ref{pf:gaussian} and the full vector/residual derivation in
Proposition~\ref{seqp:conditioning}. The finite and infinite-size
correlation domains are proved in Proposition~\ref{seqp:domain}.

Put
\[
 d_r=1+(r-1)c,\qquad d_N=1+(N-1)c.
\]
Every later recursive prediction equals
\(b_r=rca/d_r\). Consequently the fixed inclusive prediction average and
the conditional expectation of the actual mean are
\begin{equation}
 B_{r,k}=\frac{ra+kb_r}{N}
        =\frac{r d_N a}{N d_r}.
 \label{eq:seq-sim-centre}
\end{equation}
For completeness, Gaussian regression on any prefix of length \(s\)
has coefficient \(c/[1+(s-1)c]\) on the sum of that prefix. If the first
\(r\) values equal \(a\), the next prediction is \(b_r\). If the next
\(j-1\) predictions already equal \(b_r\), their prefix sum is
\(ra+(j-1)b_r=ra[d_r+(j-1)c]/d_r\). Substituting it in the regression
formula gives \(b_r\) again. Mathematical induction therefore proves the
constant-prediction assertion and \eqref{eq:seq-sim-centre}.
Proposition~\ref{seqp:recursion} supplies the regression coefficient's
matrix proof as well as the induction and Gaussian aggregate-law proof.

For an actual inclusive mean \(C_{r,k}\), the target variances are
\begin{align}
 v_{\mathrm{pop}}&=\frac{d_N}{N},
   &\operatorname{E}[C_{r,k}]&=0,
   \label{eq:seq-sim-population}\\
 v_{\mathrm{cond}}&=\frac{k(1-c)d_N}{N^2d_r},
   &\operatorname{E}[C_{r,k}\mid X_0=\cdots=X_{r-1}=a]&=B_{r,k}.
   \label{eq:seq-sim-conditional}
\end{align}
An IQ mean is obtained by applying \(x\mapsto100+15x\) to the
corresponding standardized mean; an IQ standard deviation is 15 times
the standardized SD. Variances
therefore multiply by \(15^2\), rather than by 15. The phrase
``conditional SD'' in the tables always refers to actual outcomes.
These transformations follow from the affine moment proof in
Propositions~\ref{pf:moments} and \ref{seqp:recursion}.

\section{Generating the conditional future scores directly}
The new implementation samples every finite future residual component
rather than sampling only a scalar mean with the already calculated
variance. It combines their sum with the common factor, using the
algebraically equivalent construction below without retaining a second
array of complete scores.
Let \(Z=(X_0,\ldots,X_{r-1})^{\mathsf T}\). The Gaussian conditioning
formula gives
\begin{equation}
 G\mid Z=a\mathbf1_r\ \sim\ N(\mu_r,\tau_r^2),
 \qquad \mu_r=\frac{r\sqrt c\,a}{d_r},\quad
 \tau_r^2=\frac{1-c}{d_r}.
 \label{eq:seq-sim-posterior-factor}
\end{equation}
Indeed, \(\operatorname{Cov}(G,Z)=\sqrt c\,\mathbf1_r^{\mathsf T}\)
and \(\operatorname{Cov}(Z)\mathbf1_r=d_r\mathbf1_r\). Regression
therefore gives the displayed mean and variance
\(1-rc/d_r=(1-c)/d_r\). Here \(d_r>0\) and the covariance of \(Z\)
is positive definite, so the nonsingular-block Gaussian conditioning
theorem applies \citep[Theorems 3.5.1--3.5.2]{gallager2013}.

Using fresh standard normals \(H,E_1,\ldots,E_k\), the task graph forms
\begin{equation}
 Y_j=\sqrt c\,(\mu_r+\tau_r H)+\sqrt{1-c}\,E_j,
 \qquad
 C_{r,k}=\frac{ra+\sum_{j=1}^{k}Y_j}{N}.
 \label{eq:seq-sim-future-vector}
\end{equation}
The future residuals \(E_j\) are independent of the initially observed
block, so this is the required conditional vector construction. It does
not repeatedly redraw or pretend to reobserve the fixed starting values.
In particular, it does not reject unconditioned samples until their
initial values equal 150: that event has probability zero under this
continuous model. The continuous Gaussian conditional distribution
specifies the pointwise experiment.

Each future conditional mean is \(\sqrt c\,\mu_r=b_r\). The sum of the
future random components has variance
\[
 k^2c\tau_r^2+k(1-c)
 =\frac{k(1-c)(d_r+kc)}{d_r}
 =\frac{k(1-c)d_N}{d_r}.
\]
Dividing by \(N^2\) proves \eqref{eq:seq-sim-conditional}. Similarly,
counting the \(N\) individual variance terms and \(N(N-1)\)
covariance terms in the population sum proves
\eqref{eq:seq-sim-population}. These are consequences of covariance
bilinearity, rather than results inferred from a successful simulation
\citep[Theorem 1.6.1 and Section 2.2]{durrett2019}.

\chapter{Dask implementation and the experimental design}
\label{sec:seq-sim-design}

Dask arrays express the normal draws and vector arithmetic in chunks of
\SimChunk{} observations. Dask reductions calculate counts, means and
second moments without retaining every full battery in one in-memory
object. The sample variance uses denominator \(M-1\), where \(M\)
is the number of simulated observations contributing to that estimate.
The report's tables and plotted coordinates are generated from these
reductions; PGFPlots renders those coordinates in XeLaTeX. There is no
Matplotlib calculation or plotting stage. The distinction between graph
construction and execution follows the documented Dask array and delayed
interfaces \citep{daskArray,daskDelayed}.

Each case uses \SimReplicates{} recorded seed configurations, with
\SimTrials{} observations per replicate. Results are pooled to
\SimPooledTrials{} observations per case. Pooling combines sample counts,
sums and squared sums before forming the overall variance; simply
averaging replicate SDs would not generally produce the pooled SD.
The exact within/between-sample identity establishing this fact is proved
in Proposition~\ref{sm:pooling}.
The raw replicate results remain available so the pooled estimate does
not hide replicate-to-replicate differences.

Within a replicate and starting-score count, common factor and residual
arrays are shared across the selected future counts. This supplies common
random numbers when comparing \(k=1,2,10\). The case estimates are
therefore correlated; they are not nine independent replications of one
experiment. Replicate seed configurations define separate pseudorandom
streams, whose ideal statistical interpretation assumes independent
normal observations. Neither the seed choices nor Dask scheduling proves
true stochastic independence. Versions, seed and chunk configuration are
part of the reproduction record: changing chunk boundaries can change
the pseudorandom realization while preserving the intended model.

The finite population and conditional grid contains 18 experiments, and
six more sample Gaussian laws with the analytic limiting parameters.
The total number of scalar-mean observations across that grid is
\(24\times\SimPooledTrials\). This is a count of output observations,
not of distinct people or independent comparisons between cases, and not
a count of all underlying normal variates.

\chapter{Finite-count means: the starting scores remain in the average}
\label{sec:seq-sim-means}

The following table compares observed IQ centres with
\eqref{eq:seq-sim-centre} and the population centre of 100. Its Monte Carlo
standard error quantifies uncertainty of the simulated mean estimate.
For a theoretical actual-outcome variance \(v\) on the standardized
scale, that standard error in IQ units is \(15\sqrt{v/M}\), as proved
in the common statistical-methods appendix.
In particular, the actual-outcome SD divided by this estimation SE is
\(\sqrt M\) when \(v>0\); increasing simulation size improves the
estimate without narrowing the actual-outcome distribution.


\begingroup\small
\setlength{\tabcolsep}{4pt}
\Needspace{6\baselineskip}
\noindent\textbf{Conditional centres after the observed IQ-150 block}
\begin{longtable}{rrrrrr}
\toprule
$r$ & $k$ & Target IQ & Observed IQ & MC SE (IQ) & $|z_{mean}|$ \\
\midrule
\endfirsthead
$r$ & $k$ & Target IQ & Observed IQ & MC SE (IQ) & $|z_{mean}|$ \\
\midrule
\endhead
\bottomrule
\endfoot
$1$ & $1$ & $140.0000$ & $139.9961$ & $0.0030$ & $1.2971$ \\
$1$ & $2$ & $136.6667$ & $136.6614$ & $0.0033$ & $1.5776$ \\
$1$ & $10$ & $131.8182$ & $131.8096$ & $0.0036$ & $2.3806$ \\
$2$ & $1$ & $145.8333$ & $145.8361$ & $0.0019$ & $1.4880$ \\
$2$ & $2$ & $143.7500$ & $143.7514$ & $0.0022$ & $0.6279$ \\
$2$ & $10$ & $139.5833$ & $139.5825$ & $0.0027$ & $0.3035$ \\
$3$ & $1$ & $147.7273$ & $147.7271$ & $0.0013$ & $0.1006$ \\
$3$ & $2$ & $146.3636$ & $146.3654$ & $0.0017$ & $1.0732$ \\
$3$ & $10$ & $143.0070$ & $143.0090$ & $0.0022$ & $0.8940$ \\
\end{longtable}
\endgroup

\begingroup\small
\setlength{\tabcolsep}{4pt}
\Needspace{6\baselineskip}
\noindent\textbf{Whole-battery population centres}
\begin{longtable}{rrrrrr}
\toprule
$r$ & $k$ & Target IQ & Observed IQ & MC SE (IQ) & $|z_{mean}|$ \\
\midrule
\endfirsthead
$r$ & $k$ & Target IQ & Observed IQ & MC SE (IQ) & $|z_{mean}|$ \\
\midrule
\endhead
\bottomrule
\endfoot
$1$ & $1$ & $100.0000$ & $99.9870$ & $0.0067$ & $1.9316$ \\
$1$ & $2$ & $100.0000$ & $99.9848$ & $0.0064$ & $2.3730$ \\
$1$ & $10$ & $100.0000$ & $99.9866$ & $0.0060$ & $2.2388$ \\
$2$ & $1$ & $100.0000$ & $99.9954$ & $0.0064$ & $0.7147$ \\
$2$ & $2$ & $100.0000$ & $99.9970$ & $0.0063$ & $0.4805$ \\
$2$ & $10$ & $100.0000$ & $100.0008$ & $0.0060$ & $0.1288$ \\
$3$ & $1$ & $100.0000$ & $100.0076$ & $0.0063$ & $1.2103$ \\
$3$ & $2$ & $100.0000$ & $100.0066$ & $0.0062$ & $1.0686$ \\
$3$ & $10$ & $100.0000$ & $100.0067$ & $0.0060$ & $1.1295$ \\
\end{longtable}
\endgroup



The deterministic target is above the prediction for a subsequent test
because the inclusive average still gives positive weight to the initial
150s. With one initial score and ten later scores, the target centre is
\(1450/11\), about 131.82. With two initial scores it is
\(1675/12\), about 139.58. With three it is \(20450/143\), about 143.01.
These values differ because the observed evidence differs, not because
the population mean of the standardized tests has changed.
For the strict comparison with the subsequent prediction, subtracting
\(b_r\) from \(B_{r,k}\) gives
\(ra(1-c)/(Nd_r)>0\) at the selected positive \(a\) and \(c<1\).
\Cref{seqp:exact-grid} proves the displayed fractions and target
roundings.

The largest absolute discrepancy of a pooled empirical centre from its
theoretical value is \SimMaxMeanSE{} theoretical standard errors. This
statistic is a scale-free diagnostic of the fixed recorded grid. It is
not an effect size measuring the difference between population and
conditional centres. The latter difference is deliberately introduced by
conditioning on the high initial scores. Reported numerical agreement
also does not make extra displayed decimals meaningful measures of
psychometric precision.

\chapter{Finite-count standard deviations and what they measure}
\label{sec:seq-sim-sds}


\begingroup\small
\setlength{\tabcolsep}{4pt}
\Needspace{6\baselineskip}
\noindent\textbf{Conditional SD of actual outcomes (IQ points)}
\begin{longtable}{rrrrrr}
\toprule
$r$ & $k$ & Target SD & Observed SD & 95\% lower & 95\% upper \\
\midrule
\endfirsthead
$r$ & $k$ & Target SD & Observed SD & 95\% lower & 95\% upper \\
\midrule
\endhead
\bottomrule
\endfoot
$1$ & $1$ & $6.0000$ & $6.0008$ & $5.9967$ & $6.0050$ \\
$1$ & $2$ & $6.6332$ & $6.6335$ & $6.6289$ & $6.6381$ \\
$1$ & $10$ & $7.2157$ & $7.2139$ & $7.2089$ & $7.2189$ \\
$2$ & $1$ & $3.7081$ & $3.7069$ & $3.7043$ & $3.7094$ \\
$2$ & $2$ & $4.4371$ & $4.4374$ & $4.4344$ & $4.4405$ \\
$2$ & $10$ & $5.4486$ & $5.4492$ & $5.4454$ & $5.4529$ \\
$3$ & $1$ & $2.6756$ & $2.6782$ & $2.6764$ & $2.6801$ \\
$3$ & $2$ & $3.3358$ & $3.3368$ & $3.3345$ & $3.3391$ \\
$3$ & $10$ & $4.4553$ & $4.4593$ & $4.4562$ & $4.4624$ \\
\end{longtable}
\endgroup

\begingroup\small
\setlength{\tabcolsep}{4pt}
\Needspace{6\baselineskip}
\noindent\textbf{Whole-battery population SD (IQ points)}
\begin{longtable}{rrrrrr}
\toprule
$r$ & $k$ & Target SD & Observed SD & 95\% lower & 95\% upper \\
\midrule
\endfirsthead
$r$ & $k$ & Target SD & Observed SD & 95\% lower & 95\% upper \\
\midrule
\endhead
\bottomrule
\endfoot
$1$ & $1$ & $13.4164$ & $13.4174$ & $13.4145$ & $13.4204$ \\
$1$ & $2$ & $12.8452$ & $12.8472$ & $12.8444$ & $12.8500$ \\
$1$ & $10$ & $11.9659$ & $11.9670$ & $11.9644$ & $11.9697$ \\
$2$ & $1$ & $12.8452$ & $12.8444$ & $12.8416$ & $12.8472$ \\
$2$ & $2$ & $12.5499$ & $12.5494$ & $12.5466$ & $12.5521$ \\
$2$ & $10$ & $11.9373$ & $11.9368$ & $11.9342$ & $11.9394$ \\
$3$ & $1$ & $12.5499$ & $12.5483$ & $12.5455$ & $12.5510$ \\
$3$ & $2$ & $12.3693$ & $12.3670$ & $12.3643$ & $12.3697$ \\
$3$ & $10$ & $11.9131$ & $11.9107$ & $11.9081$ & $11.9133$ \\
\end{longtable}
\endgroup



The population SD is the spread of whole-battery averages around IQ 100.
Its formula depends on total battery size, so cases with equal \(r+k\)
have the same population SD regardless of how the count is labelled.
The conditional SD is the spread around the corresponding high
conditional centre after the initial block is fixed. The two SDs refer
to different repeated-sampling experiments, not two estimators of the
same spread. The largest relative discrepancy of a pooled empirical SD
from its analytic target in the recorded grid is
\SimMaxSDRelativeError{}; the generated tables retain the individual
values, so this maximum is not a substitute for inspecting the rows.
Proposition~\ref{seqp:variance-decomposition} proves the experiment
distinction, including the variance that varying the initial block adds.

% Only Dask-generated coordinates enter these charts.  The lines join the
% three analytic finite-count values; they are not an additional numerical
% model fitted by TeX or a claim that only three possible counts exist.
\begin{figure}[htbp]
\centering
\begin{tikzpicture}
\begin{axis}[
 width=0.88\textwidth,height=5.5cm,xmode=log,
 xtick={1,2,10},xticklabels={1,2,10},
 xlabel={Subsequent scores \(k\)},ylabel={Population SD (IQ points)},
 legend columns=3,legend style={at={(0.5,1.03)},anchor=south},
 grid=major]
\addplot[blue!75!black,thick,forget plot] table[col sep=comma,x=k,y=theory_sd_iq]{simulation-data/population-r1.csv};
\addplot[blue!75!black,only marks,mark=*] table[col sep=comma,x=k,y=sd_iq]{simulation-data/population-r1.csv};
\addlegendentry{\(r=1\)}
\addplot[red!75!black,thick,forget plot] table[col sep=comma,x=k,y=theory_sd_iq]{simulation-data/population-r2.csv};
\addplot[red!75!black,only marks,mark=square*] table[col sep=comma,x=k,y=sd_iq]{simulation-data/population-r2.csv};
\addlegendentry{\(r=2\)}
\addplot[teal!80!black,thick,forget plot] table[col sep=comma,x=k,y=theory_sd_iq]{simulation-data/population-r3.csv};
\addplot[teal!80!black,only marks,mark=triangle*] table[col sep=comma,x=k,y=sd_iq]{simulation-data/population-r3.csv};
\addlegendentry{\(r=3\)}
\end{axis}
\end{tikzpicture}
\caption{Whole-battery population SD. Lines connect analytic targets;
markers are Dask estimates. The horizontal axis uses logarithmic spacing
but labels the actual future counts. All initial scores are random in
this experiment.}
\end{figure}

\begin{figure}[htbp]
\centering
\begin{tikzpicture}
\begin{axis}[
 width=0.88\textwidth,height=5.5cm,xmode=log,
 xtick={1,2,10},xticklabels={1,2,10},
 xlabel={Subsequent scores \(k\)},ylabel={Conditional SD (IQ points)},
 legend columns=3,legend style={at={(0.5,1.03)},anchor=south},
 grid=major]
\addplot[blue!75!black,thick,forget plot] table[col sep=comma,x=k,y=theory_sd_iq]{simulation-data/conditional-r1.csv};
\addplot[blue!75!black,only marks,mark=*] table[col sep=comma,x=k,y=sd_iq]{simulation-data/conditional-r1.csv};
\addlegendentry{\(r=1\)}
\addplot[red!75!black,thick,forget plot] table[col sep=comma,x=k,y=theory_sd_iq]{simulation-data/conditional-r2.csv};
\addplot[red!75!black,only marks,mark=square*] table[col sep=comma,x=k,y=sd_iq]{simulation-data/conditional-r2.csv};
\addlegendentry{\(r=2\)}
\addplot[teal!80!black,thick,forget plot] table[col sep=comma,x=k,y=theory_sd_iq]{simulation-data/conditional-r3.csv};
\addplot[teal!80!black,only marks,mark=triangle*] table[col sep=comma,x=k,y=sd_iq]{simulation-data/conditional-r3.csv};
\addlegendentry{\(r=3\)}
\end{axis}
\end{tikzpicture}
\caption{Actual-outcome SD after fixing \(r\) distinct initial scores at
IQ 150. The plotted SD concerns future randomness; the fixed list of
conditional predictions has SD zero. Lines and markers have the same
meaning as in the population plot.}
\end{figure}

For two initial scores and ten subsequent outcomes, the theoretical
population variance is \(285/2\) IQ points squared, whereas the
conditional variance is \(475/16\). Their SDs are approximately
11.937336 and 5.448624 IQ points, respectively. The actual conditional
centre is about 139.583333. A numerical estimate of that centre based on
millions of simulated batteries can be very precise while an individual
future battery remains spread around it with SD about 5.45. The fixed
list \(150,150,137.5,\ldots,137.5\) has exactly that conditional centre
as its arithmetic mean, with no Monte Carlo randomness of its own.
\Cref{seqp:exact-grid} gives exact certificates for these decimal
roundings and every analytic target in the generated tables.

The generated SD intervals are 95\% chi-square confidence intervals for
the unknown SD of an actual-outcome distribution, based on the observed
sample variance. They are not 95\% prediction intervals for individual
future means and not posterior credible intervals for a person's ability.
The exact interval formula and its Gaussian sampling assumptions are
proved in the statistical-methods appendix. Comparing many intervals
also creates a multiple-comparison issue: nominal 95\% rowwise coverage
does not imply 95\% simultaneous coverage of every row.
Corollary~\ref{sm:sd-interval} and Proposition~\ref{sm:multiple} prove precisely
the rowwise interval and multiple-row coverage statements.

\section{Why one SD decreases while the other increases}
\label{sec:seq-sim-monotonicity}
The population variance can be written
\(c+(1-c)/(r+k)\), so it strictly decreases in \(k\) at \(c=0.6\).
The conditional variance increases for the selected correlation. To see
this directly, put \(f(k)=k(d_r+ck)/(r+k)^2\), so that
\(v_{\mathrm{cond}}=(1-c)f(k)/d_r\). Expanding the finite difference gives
\[
 f(k+1)-f(k)=
 \frac{(2cr-d_r)k^2+(2cr^2+2cr-d_r)k+r^2(c+d_r)}
 {(r+k)^2(r+k+1)^2}.
\]
At \(c=3/5\), the numerator coefficients are respectively
\((3r-2)/5\), \((6r^2+3r-2)/5\) and
\(r^2(3r+5)/5\), all positive for \(r\geq1\). Thus the increase follows
from the formula, independently of the observed ordering in a simulation.
It is not a universal monotonicity claim for all admissible correlations.

Initially, the known initial block occupies a large fraction of the
inclusive mean and contributes no conditional variance. Increasing
\(k\) reduces that fixed fraction. Future measurement noise averages
away, but uncertainty in the shared factor remains. This explains how
additional future outcomes can increase the conditional spread of this
particular inclusive average even while averaging a random whole battery
reduces its population spread.

\chapter{The infinite-count cases}
\label{sec:seq-sim-limits}

Keeping \(c\) fixed while \(k\) grows requires \(0\leq c\leq1\),
as proved in Proposition~\ref{seqp:domain}. The chosen \(c=3/5\)
satisfies this condition at every finite size.

The predicted value of each subsequent score is already constant; it
does not drift with the prediction index. Increasing \(k\) instead
reduces the weight of the fixed initial block in the inclusive mean:
\[
 B_{r,k}=b_r+\frac{ra(1-c)}{d_r(r+k)}\longrightarrow b_r.
\]
The population and conditional SD limits are
\begin{equation}
 \sqrt c,\qquad \sqrt{\frac{c(1-c)}{d_r}},
 \label{eq:seq-sim-limit-sds}
\end{equation}
respectively. The first follows by letting \(N\) increase in
\(c+(1-c)/N\). For the second, rewrite the conditional variance as
\[
 \frac{1-c}{d_r}\left[
 c+\frac{1-(r+1)c}{N}-\frac{r(1-c)}{N^2}\right],
\]
and take the limit. Taking nonnegative square roots preserves these
limits, since \(\lvert\sqrt{x}-\sqrt y\rvert\leq\sqrt{|x-y|}\)
for nonnegative \(x,y\); that inequality follows by squaring and using
\((u-v)^2\leq(u-v)(u+v)\) when \(u\geq v\geq0\).

The new limit experiments sample \(\sqrt c\,H\) for the population and
\(\sqrt c(\mu_r+\tau_rH)\) after conditioning. These distributions have
exactly the limiting means and variances just derived. They do not use
an arbitrarily large finite \(k\) as a surrogate for infinity. They
also do not by themselves prove convergence: the displayed algebra
establishes the parameter limits. More explicitly, using a common
standard normal \(H\) to couple Gaussian laws with means \(m_k\) and
variances \(v_k\) gives squared mean error
\[
 \operatorname{E}\left[
 (m_k+\sqrt{v_k}H-m_\infty-\sqrt{v_\infty}H)^2\right]
 =(m_k-m_\infty)^2+(\sqrt{v_k}-\sqrt{v_\infty})^2\longrightarrow0.
\]
This proves convergence in mean square for that coupling of aggregate
laws, without asserting a particular pathwise construction for all tests.


\begingroup\small
\setlength{\tabcolsep}{4pt}
\Needspace{6\baselineskip}
\noindent\textbf{Gaussian laws with the analytically calculated limiting parameters}
\begin{longtable}{rrrrrr}
\toprule
$r$ & Law & Target IQ & Observed IQ & Target SD & Observed SD \\
\midrule
\endfirsthead
$r$ & Law & Target IQ & Observed IQ & Target SD & Observed SD \\
\midrule
\endhead
\bottomrule
\endfoot
$1$ & Cond. & $130.0000$ & $129.9906$ & $7.3485$ & $7.3449$ \\
$1$ & Pop. & $100.0000$ & $99.9852$ & $11.6190$ & $11.6133$ \\
$2$ & Cond. & $137.5000$ & $137.5006$ & $5.8095$ & $5.8087$ \\
$2$ & Pop. & $100.0000$ & $100.0011$ & $11.6190$ & $11.6174$ \\
$3$ & Cond. & $140.9091$ & $140.9116$ & $4.9543$ & $4.9601$ \\
$3$ & Pop. & $100.0000$ & $100.0060$ & $11.6190$ & $11.6324$ \\
\end{longtable}
\endgroup

% Promote the retained generated article heading without changing its data.
\begingroup
\let\subsection\section
\subsection{Specific observations in the recorded run}
With two observed IQ-150 scores and ten future scores, the simulated actual-mean centre is $139.5840$ IQ and its outcome SD is $5.4484$ IQ, compared with targets $139.5833$ and $5.4486$. These SD units describe a single actual mean, not its Monte Carlo estimation error. For the one-start limiting law, the observed centre $129.9984$ differs from target 130 by $1.408$ Monte Carlo standard errors. This small centre-estimation discrepancy is distinct from the limiting outcome SD of approximately 7.35 IQ. The largest relative SD discrepancy over all laws is $0.0287$\%; pointwise 95\% intervals are not simultaneous guarantees across all 24 laws.

\endgroup

The limiting IQ centres for \(r=1,2,3\) are 130, 137.5 and
\(1550/11\), respectively. The population SD limit is
\(\sqrt{135}\) IQ points for all three counts. Conditional limiting
variances are \(54\), \(135/4\) and \(270/11\) IQ points squared.
More distinct observed initial scores reduce remaining common-factor
uncertainty through \(\tau_r^2=(1-c)/d_r\). Equal observed values do
not mean that a single observation was simply duplicated: the model
assumes distinct tests with distinct residual terms.
Explicitly,
\(\tau_{r+1}^2-\tau_r^2=-c(1-c)/(d_rd_{r+1})<0\) for \(0<c<1\),
since \(d_{r+1}=d_r+c\). This proves the decrease in remaining factor
uncertainty at the selected correlation.

\chapter{Replication and comparison with the retained Ada run}
\label{sec:seq-sim-replication}


\begingroup\small
\setlength{\tabcolsep}{4pt}
\Needspace{6\baselineskip}
\noindent\textbf{Separate seeded runs; maxima cover all 24 laws}
\begin{longtable}{rrrrr}
\toprule
Replicate & Base seed & Max $|z_{mean}|$ & $r=3,k=10$ mean IQ & SD IQ \\
\midrule
\endfirsthead
Replicate & Base seed & Max $|z_{mean}|$ & $r=3,k=10$ mean IQ & SD IQ \\
\midrule
\endhead
\bottomrule
\endfoot
$1$ & $20261114$ & $1.445$ & $143.0099$ & $4.4548$ \\
$2$ & $20261115$ & $2.201$ & $143.0086$ & $4.4515$ \\
$3$ & $20261116$ & $1.267$ & $143.0083$ & $4.4542$ \\
$4$ & $20261117$ & $2.202$ & $143.0092$ & $4.4513$ \\
$5$ & $20261118$ & $1.859$ & $143.0083$ & $4.4542$ \\
$6$ & $20261119$ & $1.761$ & $143.0100$ & $4.4547$ \\
$7$ & $20261120$ & $1.161$ & $143.0057$ & $4.4553$ \\
$8$ & $20261121$ & $1.826$ & $143.0060$ & $4.4560$ \\
\end{longtable}
\endgroup


Replicate summaries show the direction and size of ordinary sampling
fluctuation under fixed model parameters. The pooled values use all the
observations, but replicate results are retained to reveal whether a
pooled discrepancy is broadly shared or concentrated in a stream. A
finite collection of successful seeds does not certify the generator.
The analytic mean/variance laws are the reference; no seed is selected
because it produces unusually close agreement.


\begingroup\small
\setlength{\tabcolsep}{4pt}
\Needspace{6\baselineskip}
\noindent\textbf{Retained Ada IQ examples: 200,000 observations per row}
\begin{longtable}{rrrrrr}
\toprule
$r$ & $k$ & Pop. mean & Pop. SD & Cond. mean & Cond. SD \\
\midrule
\endfirsthead
$r$ & $k$ & Pop. mean & Pop. SD & Cond. mean & Cond. SD \\
\midrule
\endhead
\bottomrule
\endfoot
$1$ & $1$ & $99.9901$ & $13.4449$ & $139.9842$ & $5.9940$ \\
$1$ & $2$ & $99.9697$ & $12.8474$ & $136.6657$ & $6.6252$ \\
$1$ & $10$ & $99.9993$ & $11.9528$ & $131.8193$ & $7.1937$ \\
$2$ & $1$ & $100.0124$ & $12.8349$ & $145.8306$ & $3.7030$ \\
$2$ & $2$ & $100.0306$ & $12.5374$ & $143.7431$ & $4.4404$ \\
$2$ & $10$ & $100.0015$ & $11.9347$ & $139.5838$ & $5.4556$ \\
$3$ & $1$ & $99.9572$ & $12.5509$ & $147.7330$ & $2.6768$ \\
$3$ & $2$ & $100.0231$ & $12.3315$ & $146.3586$ & $3.3359$ \\
$3$ & $10$ & $99.9776$ & $11.9342$ & $143.0024$ & $4.4602$ \\
\end{longtable}
\endgroup


The original Ada/SPARK campaign contains 20 finite calculation cases,
46 population/conditional sampling batches of 200,000 observations each,
and 18 identity or finite-PMF checks. It also checks independent and
perfectly correlated cases, a valid negative-correlation interior and
singular negative endpoints. The new full-vector Dask grid deliberately
concentrates on the nine requested positive-correlation IQ examples and
their limits. The earlier boundary results remain available in their
source-pinned record; this report does not reinterpret the
nonnegative-factor construction as a negative-correlation sampler.

The retained Ada program samples the scalar Gaussian aggregate law,
whereas the new finite Dask calculation draws the full future or
population residual vector and combines its sum with the common factor. Agreement between them
therefore checks two constructions of the same model. Their random-number
implementations and sample counts differ, so identical individual draws,
means or variances are not expected. The comparison is statistical and
numerical, not an exact replay of the Ada pseudorandom stream. Neither
implementation is an independent empirical dataset.

The previous native campaign discharged 430 checks across its new and
reused SPARK units, under the documented numerical and I/O contracts.
That is evidence about those recorded program obligations, not a new
formal proof of Dask, NumPy, their random generators, this report or the
psychometric assumptions. The new Python/Deal contracts, structural
inventory, static lint and runtime tests are reported separately. Their
checks cannot turn Monte Carlo agreement into universal formal proof.
Historical proof and simulation records remain immutable.

\chapter{What the results establish and what remains assumed}
\label{sec:seq-sim-interpretation}

The recorded calculations make the meaning of the three averages
concrete. Whole batteries centre at the population mean. Conditioning on
several high initial scores shifts the centre while reducing remaining
uncertainty. Filling later entries with their conditional predictions
removes that remaining outcome variation entirely from the fixed
numerical list. These statements are mathematical consequences of the
Gaussian model; the tables show that the Dask implementation reproduces
them to its reported sampling accuracy.

The experiments assume equicorrelation, Gaussian joint laws, the stated
reference scale and the specified fixed starting history. They do not
measure real-test correlation, retest effects, selection mechanisms,
ceiling validity or extreme-population frequencies. The conditional
standard deviation is not a standard error for estimating an empirical
IQ trait from collected participants. The Gaussian point-conditioning
convention and the common-factor interpretation are stated model choices.
Their appropriateness for an actual testing situation requires evidence
outside this simulation study.

\Cref{seqp:record-scope} separates the mathematical deductions
from the recorded outputs and explains why finite computational agreement
cannot establish the model's validity for real test data.

\chapter{Statistical reference laws and uncertainty calculations}
\label{sec:shared-statistical-methods}
\label{sec:simulation-statistical-methods}

These results concern ideal independent observations from the stated model.
They justify reference laws and interval calculations, rather than proving
that finite-state pseudorandom output is independent or that the model fits
real tests. The input theorems are Gaussian affine closure and the Gaussian
characteristic function \citep[Theorems 3.3.1--3.3.2]{gallager2013},
orthonormal-basis extension \citep[Theorem 6.36]{axler2024}, and linearity,
independence and elementary probability inequalities
\citep[Sections 1.6, 2.1 and 3.4]{durrett2019}. The derivations below do
not depend on a successful software test.

\begin{proposition}[Normal sample moments and the chi-square pivot]
\label{sm:normal-sample}
For an integer \(m\geq2\), independent \(Y_i\sim N(\mu,v)\), \(v>0\),
the sample mean is \(N(\mu,v/m)\). With
\(S^2=\sum_i(Y_i-\bar Y)^2/(m-1)\),
\((m-1)S^2/v\sim\chi^2_{m-1}\), independently of \(\bar Y\), and
\(\E S^2=v\).
\end{proposition}
\begin{proof}
Finite linearity and zero cross-covariances give the mean and variance;
Gaussian affine closure gives the law. Set \(Z_i=(Y_i-\mu)/\sqrt v\),
and extend the unit vector \(m^{-1/2}(1,\ldots,1)\) to an orthonormal
basis. Let \(Q\) have this vector as its first row. The transformed
Gaussian \(T=QZ\) has covariance \(QQ^{\mathsf T}=I\); its characteristic
function factors, so all its coordinates are independent standard normals.
Its first coordinate is \(\sqrt m(\bar Y-\mu)/\sqrt v\). Orthogonality
preserves squared norm, hence
\[
 \frac{1}{v}\sum_i(Y_i-\bar Y)^2
 =\sum_iZ_i^2-T_1^2=\sum_{j=2}^mT_j^2.
\]
The right side is chi-square by definition, independent of \(T_1\).
Its expectation is \(m-1\), proving unbiasedness of \(S^2\).
This argument supplies the normal and chi-square laws directly from the
referenced Gaussian and linear-algebra theorems.
\end{proof}

\begin{corollary}[SD intervals and fixed-parameter reference bands]
\label{sm:sd-interval}
Let \(q_u\) denote a \(\chi^2_{m-1}\) quantile and \(0<\alpha<1\).
An exact \(1-\alpha\) confidence interval for the unknown SD under this
normal model is
\[
 \left[\sqrt{\frac{(m-1)S^2}{q_{1-\alpha/2}}},
       \sqrt{\frac{(m-1)S^2}{q_{\alpha/2}}}\right].
\]
For a fixed theoretical \(v\), a reference interval for \(S^2\) is instead
\([vq_{\alpha/2}/(m-1),vq_{1-\alpha/2}/(m-1)]\). The reference
mean discrepancy in standard-error units is
\((\bar Y-\mu)/\sqrt{v/m}\sim N(0,1)\).
\end{corollary}
\begin{proof}
The continuous chi-square law gives
\(P(q_{\alpha/2}\leq(m-1)S^2/v\leq q_{1-\alpha/2})=1-\alpha\).
Solving these inequalities for \(v\), then taking the increasing positive
square root, gives the SD interval. Solving for \(S^2\), keeping \(v\)
fixed, gives the reference interval. The last assertion is the standardized
sample-mean law in Proposition~\ref{sm:normal-sample}. If \(v=0\), every
observation equals \(\mu\) almost surely, so this division-based pivot is
inapplicable; constant output is checked directly. SciPy evaluates the
quantiles \citep{scipystats}; the algebra, rather than the library call,
establishes the probability statement.
\end{proof}

\begin{proposition}[Bernoulli counts and sampling precision]
\label{sm:binomial}
For an integer \(m\geq1\) and \(m\) independent event indicators of common
probability \(p\in[0,1]\), the
count \(K\) is binomial, \(\E K=mp\), and
\(\Var(K/m)=p(1-p)/m\).
\end{proposition}
\begin{proof}
Independence makes the probability generating function of their sum
\((1-p+pt)^m\). The binomial theorem gives
\(P(K=j)=\binom mjp^j(1-p)^{m-j}\). Each indicator has mean \(p\)
and variance \(p(1-p)\), because its square equals itself. Summing means
and independent variances, then dividing by \(m^2\), proves the claims.
These are finite-product independence and linearity
\citep[Sections 1.6 and 2.1]{durrett2019}.
\end{proof}

\begin{proposition}[Wilson score endpoints and their coverage scope]
\label{sm:wilson}
For an integer \(m\geq1\), an observed count \(K\in\{0,\ldots,m\}\),
\(0<\alpha<1\), \(\hat p=K/m\) and
\(z=\Phi^{-1}(1-\alpha/2)\), inversion
of the normal score inequality gives endpoints
\[
 \frac{\hat p+z^2/(2m)\ \mathord\pm\
 z\sqrt{\hat p(1-\hat p)/m+z^2/(4m^2)}}{1+z^2/m}.
\]
The endpoint algebra is exact; nominal \(1-\alpha\) coverage uses a normal
score approximation and is not an exact finite-sample coverage theorem.
\end{proposition}
\begin{proof}
Squaring \(|\hat p-p|\leq z\sqrt{p(1-p)/m}\) gives
\[
 (m+z^2)p^2-(2m\hat p+z^2)p+m\hat p^2\leq0.
\]
The positive leading coefficient means precisely the interval between its
two real roots. Its discriminant is
\(4mz^2\hat p(1-\hat p)+z^4\). The quadratic formula, dividing
numerator and denominator by \(2m\), gives the stated endpoints. For
\(0<p<1\), the classical independent-identically-distributed central
limit theorem implies
\(\sqrt m(\hat p-p)/\sqrt{p(1-p)}\Rightarrow N(0,1)\), giving limiting
coverage at its continuity points, not equality for fixed \(m\).
This is the score construction of \citet{wilson1927}; the limit theorem is
\citet[Theorem 3.4.1]{durrett2019}. The implementation reuses
SciPy's binomial-result interval routine \citep{scipyBinomtest}.
\end{proof}

\begin{proposition}[Exact equal-tail binomial intervals]
\label{sm:clopper-pearson}
For an integer \(m\geq1\), a count \(j\in\{0,\ldots,m\}\), and
\(0<\alpha<1\), define \(L(0)=0\), \(U(m)=1\); for \(j>0\),
let \(L(j)\) solve \(P_{L(j)}(K\geq j)=\alpha/2\), and for \(j<m\),
let \(U(j)\) solve \(P_{U(j)}(K\leq j)=\alpha/2\). Then
\(P_p(L(K)\leq p\leq U(K))\geq1-\alpha\) for every \(p\in[0,1]\).
\end{proposition}
\begin{proof}
For \(0<j\leq m\), differentiating the finite binomial sum and
telescoping gives
\[
 \frac{d}{dp}P_p(K\geq j)
 =m\binom{m-1}{j-1}p^{j-1}(1-p)^{m-j}>0\quad(0<p<1).
\]
The tail runs continuously from zero to one, so \(L(j)\) exists uniquely;
the lower tail is strictly decreasing and gives \(U(j)\) similarly.
For \(0<j<m\), at \(p=L(j)\) the lower tail is at least
\(1-\alpha/2>\alpha/2\); its decrease implies \(U(j)>L(j)\).
The endpoint counts are ordered directly, so these really are intervals.
For a fixed \(p\), \(L(j)>p\) implies \(P_p(K\geq j)<\alpha/2\).
Such \(j\)'s form an upper terminal set, because upper tails decrease
with \(j\). If nonempty its probability is the tail at its smallest
member and is at most \(\alpha/2\). The analogous lower terminal set
\(U(j)<p\) has probability at most \(\alpha/2\). The union bound makes
total noncoverage at most \(\alpha\); endpoint parameters follow directly.
Discreteness explains why the inequality can be strict. This is the
test-inversion construction of \citet{clopper1934}, with its coverage
proof supplied here. Equivalently the endpoints are beta quantiles:
the derivative above is the \(\mathrm{Beta}(j,m-j+1)\) density and its
tail has value zero at \(p=0\), so integration gives the same cumulative
distribution; the complementary tail gives the upper endpoint.
\end{proof}

\begin{proposition}[Conditional frequencies with random denominators]
\label{sm:conditional-count}
For independent repetitions of a four-cell event experiment, condition on
\(N_A=\ell>0\) observations in event \(A\), where \(P(A)>0\).
The count also in \(B\) is then \(\mathrm{Binomial}(\ell,P(B\mid A))\).
Thus a binomial interval conditional on its observed denominator has its
stated exact or approximate coverage scope. A zero denominator gives no
defined conditional-frequency estimate.
\end{proposition}
\begin{proof}
Let \(q=P(A)\) and \(r=P(B\mid A)\). Grouping cells into
\(A\cap B,A\setminus B,A^c\), their probabilities are
\(qr,q(1-r),1-q\). Product independence and the multinomial counting
coefficient give
\[
 P(N_{AB}=j,N_A=\ell)
 =\binom m\ell\binom\ell j q^\ell(1-q)^{m-\ell}
           r^j(1-r)^{\ell-j}.
\]
Divide by \(P(N_A=\ell)=\binom m\ell q^\ell(1-q)^{m-\ell}\)
when positive. The remaining factor is the stated binomial mass. Cases
with \(q=1\) follow directly with \(\ell=m\). Conditional coverage for
each positive \(\ell\) also gives coverage averaged over positive
denominators by total probability. No convention assigning a number to
\(0/0\) supplies data when \(\ell=0\).
\end{proof}

\begin{proposition}[Pooling counts and moments]
\label{sm:pooling}
For a nonempty finite collection of groups of integer sizes \(m_j\geq2\),
means \(\bar Y_j\) and sample variances
\(S_j^2\), let \(M=\sum_jm_j\) and
\(\bar Y=M^{-1}\sum_jm_j\bar Y_j\). The pooled sample variance satisfies
\[
 (M-1)S^2=\sum_j(m_j-1)S_j^2+\sum_jm_j(\bar Y_j-\bar Y)^2.
\]
Independent equal-design Bernoulli groups pool to a binomial count of size
\(M\); independence is required for that probability law, not for the
displayed algebraic identity.
\end{proposition}
\begin{proof}
Within group \(j\), write
\(Y_{ji}-\bar Y=(Y_{ji}-\bar Y_j)+(\bar Y_j-\bar Y)\).
Squaring and summing, its cross term vanishes because
\(\sum_i(Y_{ji}-\bar Y_j)=0\). Summing groups gives the identity.
For counts, independence multiplies their generating functions to
\(\prod_j(1-p+pt)^{m_j}=(1-p+pt)^M\), proving the binomial law by
Proposition~\ref{sm:binomial}. Sharing random numbers across cases does not
establish this independence premise; in the overlapping-factor designs used
here the rows cannot be counted as additional independent replications.
Averaging group SDs generally does
not satisfy the displayed within-plus-between identity.
\end{proof}

\begin{proposition}[Simultaneous diagnostics and rare-event precision]
\label{sm:multiple}
If \(J\) fixed diagnostics have false-alarm probabilities at most
\(\alpha_j\), the probability of any false alarm is at most
\(\sum_{j=1}^J\alpha_j\), regardless of their dependence.
\end{proposition}
\begin{proof}
The pointwise bound \(1_{\cup_j F_j}\leq\sum_j1_{F_j}\), integrated,
is the union bound. Taking each \(\alpha_j=\alpha/J\) therefore gives
a family bound \(\alpha\). The same argument applies to two tails per
row, counting both in \(J\). Finite diagnostics remain observations
under a reference model and cannot establish an entire distribution.
The bound is the elementary subadditivity theorem
\citep[Section 1.1]{durrett2019}.
\end{proof}

\begin{corollary}[Zero counts, relative precision and sample-size scaling]
\label{sm:rare}
For independent Bernoulli trials, \(P(K=0)=(1-p)^m\). For \(p>0\) the
relative standard error of \(K/m\) is \(\sqrt{(1-p)/(mp)}\); requiring
it at most \(\epsilon>0\) means \(m\geq(1-p)/(p\epsilon^2)\).
At fixed \(p\in(0,1)\), multiplying \(m\) by four halves the standard
error. Zero successes do not prove \(p=0\).
\end{corollary}
\begin{proof}
The first statement is the binomial mass at zero. Divide the SD in
Proposition~\ref{sm:binomial} by \(p\), square the nonnegative inequality,
and rearrange to get the sample-size bound. The square-root dependence
proves the scaling. Every \(0<p<1\) assigns positive probability to
zero successes for finite \(m\), proving the final statement. A desired
95\% zero-count upper confidence bound, for example, follows from
\((1-U)^m=.025\) for the equal-tail construction, not from setting
\(p\) equal to zero. These conclusions concern sampling uncertainty;
they do not bound floating-point quadrature error or model error.
\end{proof}

\chapter{Complete proofs for the sequential simulation targets}
\label{sec:seq-proof-completion}

This chapter completes the deductions used in the experimental chapters.
The foundational results and their proofs appear in
\Cref{sec:proof-foundations}; the sampling and interval proofs appear
in \Cref{sec:shared-statistical-methods}. A model assumption is an
input to these deductions. A recorded numerical observation is an output
of a specified finite computation, rather than a universal theorem.

\begin{proposition}[Finite and unbounded correlation domains]
\label{seqp:domain}
For a battery of size \(N\geq2\), the equicorrelation matrix
\(R_N=(1-c)I_N+c\mathbf1_N\mathbf1_N^{\mathsf T}\) is a covariance
matrix exactly when \(-1/(N-1)\leq c\leq1\). A fixed correlation is
admissible at arbitrarily large finite sizes exactly when \(0\leq c\leq1\).
The common-factor construction in \eqref{eq:seq-sim-factor} realizes every
finite size in this latter range.
\end{proposition}
\begin{proof}
Every vector is uniquely a multiple of \(\mathbf1_N\) plus an orthogonal
vector: subtract its coordinate average from each coordinate. On the first
subspace \(R_N\) multiplies by \(d_N=1+(N-1)c\), and on the second by
\(1-c\). The spectral covariance criterion in
Proposition~\ref{pf:spectral} therefore gives exactly the two inequalities
\(d_N\geq0\) and \(1-c\geq0\), including equality. Conversely, its
square-root construction supplies a Gaussian vector even at a zero
eigenvalue. A fixed \(c<0\) violates \(d_N\geq0\) once
\(N>1-1/c\); each \(0\leq c\leq1\) satisfies both inequalities for
every finite \(N\). Independent standard normals in
\eqref{eq:seq-sim-factor} give covariance \(c\) between distinct scores
and variance one, by Proposition~\ref{pf:moments}. Thus this construction
also proves existence in the range actually simulated. It makes no claim
to sample negative correlations, even though some finite negative
correlations are admissible.
\end{proof}

\begin{proposition}[The full conditional future law and point histories]
\label{seqp:conditioning}
For \(0\leq c<1\), the Gaussian residual construction gives a conditional
distribution at every specified initial vector \(z\). At
\(z=a\mathbf1_r\), it is given by the posterior factor formula
\eqref{eq:seq-sim-posterior-factor} and the entire future-vector
construction \eqref{eq:seq-sim-future-vector}. That future vector is jointly
Gaussian with mean
\(b_r\mathbf1_k\) and covariance
\begin{equation}
 K_k=(1-c)I_k+\frac{c(1-c)}{d_r}\mathbf1_k\mathbf1_k^{\mathsf T}.
 \label{eq:seq-complete-conditional-covariance}
\end{equation}
The mean function is the unique continuous conditional-mean version.
\end{proposition}
\begin{proof}
The initial covariance \(R_r\) is positive definite: its eigenvalues
are \(1-c>0\) and \(d_r>0\). Since
\((\mathbf1_r\mathbf1_r^{\mathsf T})^2
=r\mathbf1_r\mathbf1_r^{\mathsf T}\), multiplication gives
\[
 R_r^{-1}=\frac1{1-c}
   \left(I_r-\frac c{d_r}\mathbf1_r\mathbf1_r^{\mathsf T}\right),
 \qquad R_r^{-1}\mathbf1_r=\frac1{d_r}\mathbf1_r.
\]
The coefficient of \(\mathbf1_r\mathbf1_r^{\mathsf T}\) in the product
is \(-c/d_r+c/(1-c)-c^2r/[(1-c)d_r]=0\), because
\(d_r=1-c+rc\). This also verifies the formula when \(r=1\).

In this proof \(E_1,\ldots,E_k\) denote fresh future residuals,
distinct from every initial-block residual. Set
\[
 W=G-\frac{\sqrt c}{d_r}\sum_{i<r}X_i.
\]
The covariance identities of Proposition~\ref{pf:moments} give
\(\Cov(W,Z)=\sqrt c\mathbf1_r^{\mathsf T}
-\sqrt c\mathbf1_r^{\mathsf T}R_r/d_r=0\) and
\(\Var(W)=1-rc/d_r=(1-c)/d_r=\tau_r^2\).
The vector \((W,E_1,\ldots,E_k,Z)\) is Gaussian. The future residuals
are independent of \((G,Z)\), hence their covariances with \(W,Z\)
vanish; their covariance matrix is the identity. Gaussian block
independence (Proposition~\ref{pf:gaussian}) shows that
\((W,E_1,\ldots,E_k)\) is independent of \(Z\), and that its entries
are independent with variances \(\tau_r^2,1,\ldots,1\).
Consequently the conditional law of \(G\) at \(Z=z\) is represented by
\(\sqrt c\sum_i z_i/d_r+\tau_r H\), with a fresh standard normal
\(H\) independent of all future residuals. Substitution into the original
factor equation gives the whole future conditional vector, rather than
only its separate scalar moments. Its conditional covariance is
\((1-c)\delta_{ij}+c\tau_r^2\), proving the displayed matrix.

To check conditioning explicitly, write the future vector as
\(m(Z)+V\), with \(V\) independent of \(Z\). For any bounded Borel
\(h\) and Borel initial-history set \(A\), independence and Fubini give
\[
 \E[\mathbf1_{\{Z\in A\}}h(m(Z)+V)]
 =\int_A\!\int h(m(z)+v)\,P_V(dv)\,P_Z(dz).
\]
This is the defining conditional-kernel identity, as proved generally
in Proposition~\ref{pf:conditioning}. The affine mean is continuous;
the positive Gaussian density of \(Z\) implies that two continuous
versions agreeing almost everywhere agree everywhere: a nonzero
difference would persist on an open ball of positive probability.
Individual histories nevertheless have probability zero, so division by
\(P(Z=a\mathbf1_r)\) is not a conditional-probability calculation.
The displayed kernel specifies the version evaluated there. At \(c=1\),
all scores equal \(G\) almost surely, and the compatible equal-score
history gives future scores identically \(a\) directly on that support;
no singular covariance inverse is used.
\end{proof}

\begin{proposition}[Recursive predictions and the two aggregate laws]
\label{seqp:recursion}
For fixed equal initial values, all subsequent recursively filled
predictions are \(b_r=rca/d_r\). The fixed inclusive average is
\eqref{eq:seq-sim-centre} and has SD zero. The actual population and
conditional averages are Gaussian with the means and variances in
\eqref{eq:seq-sim-population}--\eqref{eq:seq-sim-conditional}.
\end{proposition}
\begin{proof}
For any prefix length \(s\), the inverse just proved and cross-covariance
\(c\mathbf1_s^{\mathsf T}\) give next-score mean
\(c\sum_{i<s}x_i/d_s\) by Proposition~\ref{pf:conditioning}. At the
initial history this is \(rca/d_r=b_r\). If the next \(j-1\) values
are all \(b_r\), the next value is
\[
 \frac c{d_r+(j-1)c}\bigl(ra+(j-1)b_r\bigr)
 =\frac c{d_r+(j-1)c}
   \frac{ra[d_r+(j-1)c]}{d_r}=b_r.
\]
Induction proves the assertion. Averaging gives
\((ra+kb_r)/N=r(d_r+kc)a/(Nd_r)=rd_Na/(Nd_r)\).
It is a deterministic number, whose centered square is zero, proving
SD zero. At \(c=1\), the compatible-support argument proves the same
conclusion with \(b_r=a\).

For the actual population sum, covariance bilinearity gives
\(N\) variance-one terms and \(N(N-1)\) ordered distinct-pair terms,
so its variance is \(N+N(N-1)c=Nd_N\). Its mean is zero. Division
by \(N^2\) proves the population formula. For the conditional sum, the
initial block is fixed. Summing every entry of the matrix in
\eqref{eq:seq-complete-conditional-covariance} gives
\[
 k(1-c)+k^2\frac{c(1-c)}{d_r}
 =\frac{k(1-c)(d_r+kc)}{d_r}
 =\frac{k(1-c)d_N}{d_r}.
\]
Dividing by \(N^2\) and summing the means gives the conditional formula.
Both are affine transformations of the relevant Gaussian vectors, so
they are Gaussian by Proposition~\ref{pf:gaussian}. Gaussianity is
essential here: matching two moments alone would not establish a normal
law. Finally, Proposition~\ref{pf:moments} gives
\(\E[100+15C]=100+15\E[C]\) and
\(\Var(100+15C)=225\Var(C)\), proving every reporting-scale conversion.
\end{proof}

\begin{proposition}[Population spread, conditional spread and random histories]
\label{seqp:variance-decomposition}
For \(0\leq c<1\), an unconstrained random initial block gives
\begin{align*}
 \E[C_{r,k}\mid Z]&=\frac{d_N}{Nd_r}\sum_{i<r}X_i,\\
 \Var(\E[C_{r,k}\mid Z])&=\frac{r d_N^2}{N^2d_r},\\
 \Var(C_{r,k})&=\frac{r d_N^2}{N^2d_r}
                 +\frac{k(1-c)d_N}{N^2d_r}=\frac{d_N}{N}.
\end{align*}
Thus conditional and population SD describe different experiments;
conditioning does not change the stipulated marginal population SD of
each original test.
\end{proposition}
\begin{proof}
Each future conditional mean at a general \(z\) is
\(c\sum_i z_i/d_r\) by Proposition~\ref{seqp:conditioning}. The first
identity follows by adding the fixed initial sum to those \(k\) means.
There are \(r\) diagonal variance terms and \(r(r-1)\) ordered
covariance terms in the initial sum, giving variance \(rd_r\); squaring
the coefficient proves the second identity. The residual conditional
variance is constant in \(z\), by the proved covariance matrix.
Total variance (Proposition~\ref{pf:expectation}) gives the sum. Expanding
\(r d_N+k(1-c)=Nd_r\) simplifies it as displayed. At the simulated
\(c=3/5\), all factors in the between-history variance are positive,
so conditioning strictly reduces variance. The conditional law's centre
changes because the observed block is fixed high, while the unconditioned
centre remains zero. For \(r=2,k=10\), multiplication by 225 yields the
exact IQ-scale decomposition
\(285/2=1805/16+475/16\). Taking square roots gives the reported SDs;
the decomposition adds variances, not SDs.
\end{proof}

\section{Exact certificates for the numerical target columns}
\label{seqp:exact-grid}
At \(c=3/5\), \(a=10/3\), and \(N=r+k\), the preceding proposition
gives the following rational IQ-scale targets:
\[
 m=100+\frac{50r(3N+2)}{N(3r+2)},\qquad
 V_P=\frac{45(3N+2)}N,\qquad
 V_C=\frac{90k(3N+2)}{N^2(3r+2)}.
\]
Indeed, \(d_j=(3j+2)/5\); inserting this into the proved centre and
variance formulas and multiplying variances by 225 gives the three
expressions. Direct substitution gives all finite exact entries below.
For \(k=\infty\), the proved limits in \Cref{sec:seq-sim-limits}
give \(m=100+150r/(3r+2)\), \(V_P=135\), and
\(V_C=270/(3r+2)\).

\begin{center}\small
\begin{tabular}{rr r r r}
\toprule
\(r\)&\(k\)&\(m\)&\(V_P\)&\(V_C\)\\\midrule
1&1&\(140\)&\(180\)&\(36\)\\
1&2&\(410/3\)&\(165\)&\(44\)\\
1&10&\(1450/11\)&\(1575/11\)&\(6300/121\)\\
2&1&\(875/6\)&\(165\)&\(55/4\)\\
2&2&\(575/4\)&\(315/2\)&\(315/16\)\\
2&10&\(1675/12\)&\(285/2\)&\(475/16\)\\
3&1&\(1625/11\)&\(315/2\)&\(315/44\)\\
3&2&\(1610/11\)&\(153\)&\(612/55\)\\
3&10&\(20450/143\)&\(1845/13\)&\(36900/1859\)\\\midrule
1&\(\infty\)&\(130\)&\(135\)&\(54\)\\
2&\(\infty\)&\(275/2\)&\(135\)&\(135/4\)\\
3&\(\infty\)&\(1550/11\)&\(135\)&\(270/11\)\\\bottomrule
\end{tabular}
\end{center}

Here is an exact rounding proof, independent of floating-point evaluation.
For \(m=u/v\), \(v>0\), a displayed \(t/T\) differs from \(m\)
by less than half a unit precisely when
\(D=v-2|Tu-tv|>0\). For an SD \(\sqrt{p/q}\), \(p,q>0\), and
positive displayed integer \(t\), the two certificates
\[
 L=4T^2p-(2t-1)^2q>0,\qquad
 U=(2t+1)^2q-4T^2p>0
\]
are equivalent to
\((t-1/2)/T<\sqrt{p/q}<(t+1/2)/T\): multiply by \(4T^2q\)
and square the positive endpoints. Thus positivity proves the nearest
rounding, with no tie convention needed. The following exact integers
certify all four-decimal target mean and SD entries in the tables, using
\(T=10^4\), the fractions above, and their displayed integers \(t\).
The population mean is exactly 100 in every row.

\begin{center}\scriptsize
\setlength{\tabcolsep}{4pt}
\begin{tabular}{rr r rr rr}
\toprule
\(r\)&\(k\)&\(D\)&\(L_P\)&\(U_P\)&\(L_C\)&\(U_C\)\\\midrule
1&1&1&621071&452241&239999&240001\\
1&2&1&848591&179025&528431&2225\\
1&10&7&1108621&9421371&24721751&45126225\\
2&1&2&848591&179025&584316&602276\\
2&2&4&1011982&996002&558704&5120784\\
2&10&4&1649950&260018&5138544&1835664\\
3&1&5&1011982&996002&9206676&211436\\
3&2&3&661775&327769&982625&13694895\\
3&10&123&12046427&343197&95794525&566797691\\
1&\(\infty\)&1&359&929161&113039&474841\\
2&\(\infty\)&2&359&929161&465116&1393924\\
3&\(\infty\)&9&359&929161&3790525&569259\\\bottomrule
\end{tabular}
\end{center}
Every listed integer is positive; the established equivalences therefore
prove the claimed roundings. For the six-decimal worked example in
\Cref{sec:seq-sim-sds}, the corresponding SD pairs \((L,U)\),
with \(T=10^6\), are \((169283518,21713858)\) for
\(\sqrt{285/2}\approx11.937336\) to the stated precision and
\((125135856,572288016)\) for \(\sqrt{475/16}\approx5.448624\).
For \(1675/12\) displayed as \(139.583333\), the mean certificate is
\(D=4\). These are strict rounding statements; the decimals do not equal
the irrational SDs exactly.

The theoretical MC SE columns use \(\sqrt{V/M}\) with
\(M=4\,000\,000\), so the same certificate applies with denominator
\(Mq\). In finite-grid row order above, their four-decimal display
integers \(t\) are
\[
 \begin{aligned}
  \text{Population:}\quad &(67,64,60,64,63,60,63,62,60),\\
  \text{Conditional:}\quad &(30,33,36,19,22,27,13,17,22).
 \end{aligned}
\]
Substitution of the listed exact variances into the two integer expressions gives
\(\min L=96\,000\,000\) and \(\min U=356\,000\,000\) across
these 18 entries. Both minima are positive, which proves every displayed
MC SE rounding by the same inequalities. Observed estimates, the
standardized discrepancies calculated from them, and confidence-interval
endpoints calculated from observed variances remain finite numerical
records; their reference-law interpretation is proved separately.

\section{What a proof of a recorded calculation means}
\label{seqp:record-scope}
The target columns and their count/limit comparisons are mathematical
deductions above. The observed columns, extrema, replicate comparisons,
CSV coordinates, graph sizes, seeds, software versions and historical
proof-check counts are finite reproduction records. Their evidential
claim is that the named source, inputs and reduction formulas produced
the recorded outputs; it is checked by reproducing that calculation and
comparing the record. A theorem about ideal Gaussian samples does not
deduce a particular seeded output. Conversely, agreement of finitely many
outputs cannot prove an identity for all model inputs. For example, a
function altered only at an untested input has exactly the same tested
outputs. This counterexample explains the logical scope of finite tests.

There are \(3\cdot3\cdot2=18\) finite laws and \(3\cdot2=6\)
limiting laws, hence 24 output cases; four one-million-observation
replicates give four million per case and 96 million scalar-mean output
observations. Reusing normal arrays across cases does not alter that
output count, but does prevent interpreting the cases as independent
replications. Pooling formulas are proved in
Proposition~\ref{sm:pooling}; ideal sample laws, interval inversion and
the scope of multiple comparisons are proved in
Proposition~\ref{sm:normal-sample}, Corollary~\ref{sm:sd-interval} and
Proposition~\ref{sm:multiple}. No event-count denominator is used in this sequential
experiment: every generated scalar mean contributes to its row. The
probability-zero initial history is handled by the conditional kernel,
not by accepting or rejecting population observations. Automated software
obligation counts concern the recorded source and its contracts; they
are neither proofs of this written derivation nor empirical validation of
the Gaussian model.

\section{Map from experimental claims to proofs and evidence}
\label{sec:seq-sim-proof-map}

Every mathematical claim in the experimental exposition belongs to a row
below. The numbered proofs are part of this standalone report; consulting
the companion analysis is optional. The primary references identify the
well-known results from which the proofs proceed. Definitions, model
choices and finite observed records have their own evidential status.

\begingroup\small
\setlength{\tabcolsep}{4pt}
\begin{longtable}{>{\raggedright\arraybackslash}p{.25\textwidth}
                 >{\raggedright\arraybackslash}p{.38\textwidth}
                 >{\raggedright\arraybackslash}p{.28\textwidth}}
\toprule
Exposition/claim & Proof or exact record classification & Established foundation\\
\midrule\endfirsthead
\toprule
Exposition/claim & Proof or exact record classification & Established foundation\\
\midrule\endhead
\bottomrule\endfoot
Abstract and Section~\ref{sec:seq-sim-questions}: standardization, fixed
starting values, three means and IQ units &
Assumptions/definitions are stated there. Their consequences are proved in
Propositions~\ref{seqp:recursion} and \ref{seqp:variance-decomposition};
the exact grid is in Section~\ref{seqp:exact-grid}. &
Moment identities and Gaussian transformations,
\ref{pf:moments}, \ref{pf:gaussian}
\citep{durrett2019,gallager2013}.\\\midrule
Section~\ref{sec:seq-sim-model}: factor construction, admissible finite
correlations and the fixed-correlation infinite domain &
Proposition~\ref{seqp:domain} proves the two eigenspaces, the PSD
inequalities and the nonnegative factor realization. &
Spectral theorem/positive square root, \ref{pf:spectral}
\citep{axler2024}; covariance expansion.\\\midrule
Section~\ref{sec:seq-sim-model}: posterior factor, whole future block,
independence and the probability-zero initial history &
Proposition~\ref{seqp:conditioning} proves the inverse, residual
independence, full conditional kernel, covariance matrix and continuous
mean-version uniqueness; it treats \(c=1\) directly on its support. &
Gaussian characteristic functions, regression residual and conditional
integral identity, \ref{pf:gaussian}, \ref{pf:conditioning}
\citep{gallager2013,durrett2019}.\\\midrule
Sections~\ref{sec:seq-sim-model}, \ref{sec:seq-sim-means}: every prediction,
inclusive centre, population centre, normality and SD zero &
Proposition~\ref{seqp:recursion} proves regression plus induction, the
inclusive numerator and denominator, Gaussian closure and constant-list
variance. The positive difference from the next prediction is computed
explicitly in Section~\ref{sec:seq-sim-means}. &
Conditional regression and moment identities,
\ref{pf:conditioning}, \ref{pf:moments}; elementary induction.\\\midrule
Sections~\ref{sec:seq-sim-model}, \ref{sec:seq-sim-sds}: actual population
and conditional mean SDs; dependence on total count &
Proposition~\ref{seqp:recursion} counts all diagonal and ordered
off-diagonal covariances; Proposition~\ref{seqp:variance-decomposition}
separates between-history and remaining variance. &
Covariance bilinearity, conditional expectation and total variance,
\ref{pf:moments}, \ref{pf:expectation}
\citep{durrett2019}.\\\midrule
Section~\ref{sec:seq-sim-monotonicity}: decreasing population and
increasing conditional SD at \(c=3/5\) &
The displayed finite-difference identity and its three strictly positive
coefficients prove the conditional assertion for all integer \(r,k\geq1\).
The population difference is negative from \(c+(1-c)/(r+k)\). &
Proved variance formulas; order-preservation of the nonnegative square root.
No inference from plotted ordering is required.\\\midrule
Section~\ref{sec:seq-sim-limits}: prediction and inclusive-mean limits,
both SD limits, limiting Gaussian experiments and common-factor uncertainty &
That section proves the exact \(1/N\) error, conditional variance
expansion, square-root inequality and the common-normal mean-square
coupling. Proposition~\ref{seqp:domain} establishes the permissible
fixed-correlation range. &
The proved finite targets, Gaussian affine closure and moment identities,
\ref{pf:gaussian}, \ref{pf:moments}.\\\midrule
Sections~\ref{sec:seq-sim-means}--\ref{sec:seq-sim-limits}: all rational
targets and displayed analytic decimal approximations &
Section~\ref{seqp:exact-grid} proves substitutions and supplies exact
positive-integer certificates for every table target rounding and the
six-decimal worked example. &
Affine variance scaling and the proved target formulas; rational
arithmetic and order of positive squares.\\\midrule
Sections~\ref{sec:seq-sim-design}--\ref{sec:seq-sim-sds}: sample SD,
theoretical MC SE, chi-square SD intervals and their scope &
\ref{sm:normal-sample} proves the normal/chi-square sampling laws;
\ref{sm:sd-interval} inverts the pivot. \ref{sm:multiple} proves the
multiple-comparison bounds without assuming independent case estimates. &
Orthogonal basis extension \citep{axler2024}, Gaussian block independence
\citep{gallager2013}, quantile inversion and the union bound.\\\midrule
Sections~\ref{sec:seq-sim-design}, \ref{sec:seq-sim-replication}: pooling,
four replicates, output counts and shared random numbers &
\ref{sm:pooling} proves count/sum/within-between-square identities;
Section~\ref{seqp:record-scope} proves the finite grid arithmetic. The
seed/chunk/sharing choices are implementation records. &
Finite-sum identities. Independent-observation sampling laws are stated
ideal assumptions, rather than consequences of seed values.\\\midrule
Tables, figures, generated discussion and retained Ada comparison &
Observed estimates, maxima, intervals, CSV coordinates, seeds and proof
counts are finite reproduction records. Their exact scope and the
counterexample to inferring universal correctness from tests appear in
Section~\ref{seqp:record-scope}. &
Documented Dask/SciPy interfaces \citep{daskArray,daskDelayed,scipystats};
source-pinned program evidence has its stated separate trust boundaries.\\\midrule
Section~\ref{sec:seq-sim-interpretation}: what the experiment establishes,
model adequacy and distinction between predictions and outcomes &
Propositions~\ref{seqp:recursion} and \ref{seqp:variance-decomposition}
prove the model conclusions. Section~\ref{seqp:record-scope} separates
finite observations from universal claims and empirical validation. &
The stated Gaussian model and proved conditional laws. No theorem here
asserts that real IQ data satisfy the assumptions.\\
\end{longtable}
\endgroup

The shared probability-count results \ref{sm:binomial}, \ref{sm:wilson},
\ref{sm:clopper-pearson}, \ref{sm:conditional-count} and \ref{sm:rare}
are proved for completeness and use in the other two simulation reports.
This sequential report uses means and SDs of all generated outcomes,
rather than accepted-event proportions or random conditional-count
denominators. Its interval claims use the Gaussian sample results above.


\bibliographystyle{plainnat}
\bibliography{simulation_references}

\appendix
\chapter{AI disclosure}
\label{sec:ai-disclosure}
This report, its Dask implementation and its review were prepared with
OpenAI Codex and collaborating AI agents. Anders Hellström specified the
questions and required separate simulation reports. No new participant
data were collected. Independent human peer review and empirical
psychometric validation are not claimed.

% Compatibility alias: canonical AI-use disclosure is in library/latex.
% Canonical full appendix text; the chapter heading is defined in report_style.tex.
This report and its supporting repository were substantially generated and
revised with OpenAI Codex and collaborating AI agents. The session identifies
Codex as based on the GPT-6 family. Exact serving checkpoints, per-agent model
identifiers and generation settings were not preserved; no more precise model
attribution is claimed. This disclosure covers the mathematical exposition,
simulations, programs, documentation, verification, typesetting and publication,
including subsequent layout and documentation revisions.

AI contributions included interpreting the problem and historical articles;
developing and reviewing derivations, counterexamples, theorem references,
proofs and claim maps; calculating examples and limits; editing
Python/Dask/Deal implementations, contracts, tests and auditors; executing
the recorded checks; producing tables and scientific charts; preparing the
XeLaTeX/KOMA-Script books; visually reviewing PDFs; and preparing repository
updates.

The requester Anders Hellström supplied questions and sources, clarified notation and
averaging, selected example parameters, specified formats and verification
requirements, requested revisions and authorized publication. Repository
ownership and commit attribution do not identify every passage as human-authored
or independently verified. Parallel mathematical and visual reviews were
performed by AI agents and automated tools within the same workflow.
Independent human peer review, professional sign-off and endorsement by
original authors, source hosts or cited organizations are not recorded and
are not claimed.

Original authors and references remain credited. Historical articles are
analysed rather than relicensed; replacement articles are AI-generated
analytical reconstructions. Sources were consulted at different levels,
including mathematical texts, articles, official software documentation and
bibliographic metadata. Reference notes identify versions and access limits;
full-text examination of every cited work is not claimed. Mathematical proofs
establish consequences of stated hypotheses, rather than empirical adequacy
of the IQ model.

Verification claims rest on retained commands, actual outputs, tool versions,
source hashes and checksums. Previous disconnected implementations and
their proof campaigns were removed from maintained source; they remain accessible
through earlier Git history and are not current-source verification for
Python/Dask. Deal contract declarations, static lint, runtime/property
tests, exact Rocq kernel checks, numerical accuracy, PRNG properties and
model validity have separate scopes. Unsupported, skipped or unrun checks
are not reported as successful proofs. Trusted numerical/library boundaries
and specification adequacy remain explicit.
Independent numerical validation means a separate implementation, algorithm or
library, not independent human authorship or review.

No participant data were collected and no empirical psychometric validation
study was conducted.
Simulations use synthetic observations under stated probability models. A typesetting-only revision does not, by itself, execute additional random
simulations or prove additional Rocq theorems. Older
compiler evidence is available through Git history, not active source files.
The repository does not preserve a complete prompt/tool transcript, exact per-line authorship
or AI-contribution percentage; these have not been invented. Its license
applies to original material, subject to third-party rights.


\end{document}
